\chapter{Síntesis: la estructura discreta del continuo}
\label{chap:sintesis-hmt}

\section{Un núcleo finito con memoria}

La Holografía Modular Triádica parte de tres objetos definidos antes de
cualquier aplicación. APP aporta el toro aritmético nonádico y su categoría
de caminos cardinales. El TRIT registra la orientación local sin identificar
la firma visible con el estado completo. El TPK transporta posición, hoja
aritmética, residuo, cociente, arrastre, orientación, fase, frontera,
supervivencia y registro dodecafásico. La composición
\[
\mathrm{APP}\longrightarrow\mathrm{TRIT}
\longrightarrow\mathrm{TPK}
\]
no abrevia tres nombres para un mismo artefacto: distingue soporte,
orientación y dinámica.

Sobre una misma posición, las evaluaciones aditiva y multiplicativa generan
particiones diferentes. La suma actúa por traslaciones; el producto pliega
las fibras del radical del anillo local. La extensión
\[
n=R+9Q
\]
conserva el residuo visible y el cociente euclídeo. El cociclo de arrastre
garantiza la composición entre celdas, mientras que la reducción módulo tres
extrae la orientación y el lector módulo mil agrupa la expansión
arquimediana. De este modo, módulo nueve, módulo tres y base mil ocupan
funciones matemáticas distintas dentro de una sola dinámica.

Los generadores norte, sur, este y oeste, los dos cursores aritméticos y los
ciclos de veintisiete y ciento ocho pasos convierten la célula local en una
teoría de transporte. El emisor finito recorre
\[
104\,976\longrightarrow468\longrightarrow243
\]
estados, emisiones decimales y palabras ternarias, respectivamente. La
factorización \(468=18\cdot26\) separa las órbitas aditivas de las firmas
multiplicativas. El retorno nonádico completa nueve fases: vuelve la fase,
pero la frontera conserva la historia. El retorno fibrado distingue retorno
de fase y retorno de estado: el transporte acumulado es un endomorfismo
vertical y constituye monodromía cuando es invertible.

\section{Número, continuidad y clasificación aritmética}

Un número HMT es una sección compatible del sistema de supervivencia. El
lector y la genealogía transportada son datos tipados que acompañan su
evaluación, pero no alteran la definición de la sección. Sus cilindros
anidados determinan un valor arquimediano cuando son no vacíos, están
anidados y su diámetro tiende a cero. El valor es la imagen de evaluación de
la sección; la sección retiene además las coordenadas que la evaluación
olvida. En la carta Hensel bruta se demuestra la cobertura de cada intervalo
fundamental. Para el sistema superviviente calibrado, la formulación de la
recta real como cociente arquimediano queda condicionada a la cobertura,
prolongabilidad y compatibilidad a toda profundidad.

La clasificación ordinaria de los números se conserva y se refina.
Terminación, periodicidad, algebraicidad y trascendencia continúan
determinadas por sus criterios aritméticos. HMT añade periodicidad del estado,
monodromía de fase, multiplicidad de la fibra, supervivencia y memoria de
frontera. Una expansión periódica puede proceder de un estado monodrómico; un
retorno local de palabra puede coexistir con una genealogía aperiódica; dos
secciones distintas pueden compartir valor sin compartir ruta.

La completación
\[
1000=729+271=729+243+27+1
\]
expresa la conservación normalizada de la unidad entre interior, estratos de
frontera y cierre. Su extensión binomial
\[
10^d=(9+1)^d
\]
organiza el mismo principio en rango arbitrario. La sucesión no escindida
\[
0\longrightarrow\mathbb F_3\longrightarrow
\mathbb Z/9\mathbb Z\longrightarrow\mathbb F_3\longrightarrow0
\]
demuestra que la multiplicación requiere la coordenada de cociente. Dentro
de la hipótesis estructural de incidencia declarada, los invariantes \(54\) y \(1215\)
seleccionan de manera única el rango seis mediante
\[
9d=54,\qquad81\binom d2=1215.
\]

\section{Genealogía conjunta de las constantes}

En el perfil TPK enriquecido, las regiones nombradas de \(\pi\), \(e\) y
\(\varphi\), junto con sus semillas, son primitivas declaradas con control
externo. La relación EF--G9 define su prolongación por elevación, cilindro,
frontera, supervivencia y evaluación, y los certificados prueban esa
operación en las ventanas finitas publicadas. La microfibra de \(\pi\)
contiene cinco regiones distintas sobre la misma palabra visible, con
partición
\[
3_{++}+1_{--}+1_{\perp},
\qquad
432+4\cdot144=1008.
\]
La región de \(e\) codifica propagación y exhibe un cuadrado local de palabra
seguido por ruptura de memoria; la de \(\varphi\) codifica la autoescala. La
conexión nonádica proporciona la regla coinductiva de prolongación. Si los
niveles estables permanecen no vacíos y compatibles a toda profundidad,
determina las secciones globales; la verificación computacional vigente
certifica ventanas finitas y no sustituye esa hipótesis global.

El cierre dodecafásico transforma el estado firmado mediante tres bloques de
Hadamard y obtiene
\[
K=(234,543,140,729,659,824,621,58,914,794,146,601).
\]
Las cartas
\[
U_{12}^{\pm}
\longleftrightarrow K
\longleftrightarrow(D_3K,D_4K,Q)
\]
son reversibles sobre los datos declarados. La recurrencia de arrastre en
base mil integra después las lecturas regionales y la condición de frontera,
y produce el racional finito
\[
\alpha_{12}
=0.007297352569283800997285105472380663,
\]
\[
\alpha_{12}^{-1}=137.035999084\ldots.
\]
Ésta es una identidad exacta relativa al estado dodecafásico, las regiones y
la frontera declarados. La coincidencia de la cadena inicial de
\(\alpha_{12}^{-1}\) con el valor central de CODATA 2018 es un
\textsc{reconocimiento externo}; no constituye una predicción metrológica
independiente ni una prueba sin valor objetivo de la procedencia de las primitivas.

Relativamente a esas entradas y a la regla de acarreo, \(\alpha_{12}\) es el
invariante de compatibilidad entre clausura, propagación, autoescala y memoria
dodecafásica. La relación bidireccional
\[
\mathscr R(\pi,\varphi)=\frac{\pi}{\varphi^2},
\qquad
(\mathscr R\circ\mathsf J)(\pi,\varphi)=\frac{\varphi}{\pi^2}
\]
enlaza el retorno areal--volumétrico y el exponente de la semilla
electrónica. La identidad
\[
\frac{\varphi}{\pi^2}
=\frac{1}{\varphi^3(\pi/\varphi^2)^2}
\]
expresa las dos orientaciones de un mismo lector, no dos coincidencias
separadas.

\section{La segunda evaluación de la frontera}

La lectura arquimediana contrae cilindros compatibles. La lectura excepcional
conserva incidencia. En el dominio calibrado ambas reciben el mismo estado
enriquecido y se separan antes de la proyección ternaria visible. El teorema
de obstrucción demuestra que una aplicación que factorice por la palabra
visible posee exactamente las fibras que esa palabra ya había identificado;
por ello no puede reconstruir después la información excepcional. El testigo
común demostrado es finito; la identificación de un dominio común a toda
profundidad conserva el estatuto de teorema condicional.

Relativamente al calibre de Paley--Witt declarado, la construcción produce un
operador \(J\) con
\[
J^2=-I.
\]
La raíz interna de menos uno aparece así como rotación del módulo. La
reducción ternaria construye el código extendido de Golay y el diseño
\(W_{12}\); las involuciones de estrella generan \(M_{12}\). Desde el código
común se abren dos prolongaciones:
\[
C_W\longrightarrow W_{12}\longrightarrow W_{24},
\]
\[
C_W\longrightarrow N(A_2^{12})
\longrightarrow\Lambda_{24}
\longrightarrow V^\natural
\longrightarrow\mathbb M.
\]
La primera conserva la incidencia de hexadas y octadas. La segunda construye
un retículo de Niemeier, un vecino de índice tres par, unimodular y sin
raíces, el retículo de Leech y, por los teoremas clásicos correspondientes,
el módulo vértice lunar y el Monstruo.

La bandera
\[
B_K\subset H_\alpha^-
\]
es el testigo finito común. El transporte positivo de proyectores amplía ese
testigo mediante
\[
P_G\longleftrightarrow P_3\longleftrightarrow Q_3,
\qquad
\operatorname{tr}_{11}(P_G)
=\operatorname{tr}_{11}(P_3)
=\operatorname{tr}_{11}(Q_3)
=\frac3{11}.
\]
La compatibilidad entre regiones, dodecafase e incidencia deja de reducirse
así a una igualdad de cardinales: se expresa por isometrías parciales y
proyectores de rango tres.

El mismo sector contiene cinco realizaciones exactas de \(-1/12\) en
dominios distintos. La quinta utiliza el ciclo dodecafásico: si \(S\) es el
desplazamiento de orden doce, los proyectores sobre los invariantes de
\(S^3\) y \(S^4\) tienen rangos tres y cuatro, de modo que
\[
\frac1{12}\operatorname{Tr}
\bigl(\Pi_{\ker(I-S^3)}-\Pi_{\ker(I-S^4)}\bigr)
=\frac{3-4}{12}=-\frac1{12}.
\]
La independencia de esta demostración reside en su álgebra racional de
proyectores.

\section{Corredor constructivo e índice común cincuenta y cuatro}
\label{sec:coherencia-tipada-defecto-54}

La rama arquimediana construye \(\pi\), \(e\), \(\varphi\) y el cierre
dodecafásico de \(\alpha_{\mathrm{HMT}}\). Sobre el mismo dominio
enriquecido, la rama excepcional conserva incidencia y transporta el
generador ternario desde APP y TPK hasta el vecino HMT, el orbifold y
Moonshine. El funtor ciclotómico no sustituye los lectores arquimedianos:
refuerza constructivamente la segunda evaluación y hace explícita su
genealogía micro--macro.

\begin{corollary}[Índice escalar común del cierre micro--macro]
\label{prop:coincidencia-escalar-tipada-54}
Conservando los dominios y las estructuras de partida de las cuatro
construcciones, se cumple
\[
\boxed{
 D_{\rm APP}
 =[q]\,t_3(q)
 =v_\alpha^2
 =\frac{196884}{5\cdot729+1}
 =54,}
\]
donde
\[
 D_{\rm APP}
 :=\sum_{i,j=1}^{9}\operatorname{dr}_9(ij)
   -\sum_{i,j=1}^{9}\operatorname{dr}_9(i+j),
\]
\(t_3(q)=(\eta(\tau)/\eta(3\tau))^{12}\) y \([q]\) extrae el
coeficiente de \(q\).
\end{corollary}

\begin{proof}
El censo de APP del \cref{chap:app} da
\(D_{\rm APP}=459-405=54\). La expansión modular del
\cref{chap:nthw} da \([q]\,t_3(q)=54\). El
\cref{sec:vecino-leech-ejecutable} demuestra
\(v_\alpha^2=54\), y el \cref{thm:corredor-hmt-moonshine} conserva
\[
 196884=54(5\cdot729+1).
\]
La composición de estas cuatro igualdades prueba la afirmación escalar.
\end{proof}

El \cref{thm:compatibilidad-rectora-micro-macro} construye la compatibilidad
composicional en los dominios representacionales, de Fock, reticulares y
modulares. Sus tramos construidos son: el transporte fase--longitud sobre
caminos TPK; el acoplamiento APP--Fock agregado; la subdivisión de cada
macrotick en cuatro microticks; y el alineamiento APP--Witt que preserva el
vecino reticular. El corolario conserva después las igualdades escalares
tipadas, sin postular morfismos entre sus codominios:
\[
\begin{array}{c|c}
\text{objeto y residencia}&\text{valor demostrado}\\ \hline
D_{\rm APP}&54\\
\operatorname{Tr}_{o,\mathcal F_2}&54\\
\lbrack q\rbrack\,t_3&54\\
v_\alpha^2&54\\
196884&54+10\cdot3^9
\end{array}
\]
El no descenso del defecto APP local por el cociente fase--longitud permanece
como control exacto: la compatibilidad agregada no convierte ese defecto en
una función del estado cociente ni en un escalar que viaje sin mapa entre
Fock, retículos y formas modulares.

\begin{remark}[Normalización \(744\)]
El \cref{chap:canales-enteros-constantes} prueba
\[
 C_\pi+C_e+C_\varphi-3
 =315+271+161-3
 =744
 =729+15,
\]
mientras el \cref{thm:corredor-hmt-moonshine} utiliza la normalización
\[
 J_{\rm Mon}=j-744.
\]
La coincidencia de los enteros es exacta y se conserva como referencia
cruzada. Convertirla en una identidad causal exigiría un mapa adicional que
transportase las tres unidades apuntadas al término constante del
Hauptmodul.
\end{remark}

\begin{remark}[Suma de alternativas, producto de transportes y gravedad]
La linealización de la categoría de caminos realiza exactamente dos
operaciones: la suma reúne historias alternativas y la concatenación
multiplica sus pesos o transportes. Esta estructura algebraica sustenta la
construcción de amplitudes ya demostrada en su residencia principal. No
identifica, por sí sola, tales operaciones con las hojas aritméticas de APP
ni con las composiciones gravitatorias. Para establecer esa identificación
debe construirse un funtor compatible, o un homomorfismo de semianillos, que
preserve suma, producto, unidades y las equivalencias declaradas. Mientras
falte ese mapa, el puente suma/producto--gravedad es una analogía algebraica
precisa, no un teorema de realización.
\end{remark}

\section{Rango, dualidad y frontera dimensional}

Las pantallas \(12\to11\to10\) y \(26\to24\) son aplicaciones tipadas. La
primera elimina el modo uniforme y polariza después un sector mediante
\(K\). La segunda realiza el cociente isotrópico
\[
e_+^\perp/\mathbb Ze_+\cong\Lambda_{24}
\]
en \(II_{25,1}\). La involución
\[
\mathsf T(m,w,R)=(w,m,R^{-1})
\]
intercambia coordenada visible y memoria y conserva la forma espectral de
radio. Estas tres construcciones proporcionan la interfaz algebraica con las
reducciones de rango y las dualidades de la teoría M.

La frontera con Mecánica Dimensional se alcanza después del cierre relativo
de \(\alpha_{12}\). En la notación usada por los lectores posteriores,
\(\alpha_{\mathrm{HMT}}:=\alpha_{12}\), y la escala decimal define
\[
A=1000\alpha_{\mathrm{HMT}}.
\]
El selector regional, el determinante de retorno y la recurrencia fijan,
relativamente a las reglas del lector analítico graduado, la fase
complementaria primero en radianes y después en grados. El lector
determinantal local declarado produce
\[
\Sigma_H=\varphi\alpha_{\mathrm{HMT}}^{16}
\]
y su década \(-34\) antes de la carta de unidades. Quedan así disponibles dos
entradas matemáticas tipadas para Mecánica Dimensional: la interfaz angular y
la recta abstracta de acción. Su identificación con observables y unidades
físicos pertenece, respectivamente, a una \textsc{interfaz física} y a una
\textsc{transducción metrológica} posteriores.

\section{Conclusión}

La obra reúne resultados de estatutos distintos dentro de una arquitectura
común. Son \textsc{exactos internos} los censos finitos, las identidades
algebraicas y las construcciones de incidencia una vez fijados sus dominios.
El cierre dodecafásico, el contraángulo y el orden determinantal son exactos y
se demuestran a partir de una \textsc{estructura de partida explícita}, con
lectores declarados. La comparación con
CODATA es \textsc{reconocimiento externo}; la asignación de unidades es
\textsc{transducción metrológica}; y el paso hacia partículas, observables y
campos es una \textsc{interfaz física}.

El cierre del infinito se formula coinductivamente: ningún nivel finito agota
la expansión, y la relación de transición especifica cómo debe prolongarse
cada estado. Las ventanas certificadas prueban la operación local. La
existencia, cobertura o unicidad de secciones a toda profundidad se obtiene
cuando se verifican las hipótesis globales del sistema inverso. Promover las
interfaces dimensionales a una realización física completa permanece como
\textsc{programa abierto}.
